\section{数据结构}

\subsection{线段树整合}

\begin{lstlisting}[language=C++]
struct Tag {
    int x;
    void apply(const Tag &t) {}
};

struct Info {
    int l, r;
    void apply(const Tag &t) {}
};
Info operator+(const Info &u, const Info &v) {
    Info res;
    res.l = u.l;
    res.r = v.r;
    return res;
}

struct LazySegmentTree {
    int n;
    vector<Info> info;
    vector<Tag> tag;
    LazySegmentTree(int n_, Info v_ = Info()) {
        init(n_, v_);
    }
    LazySegmentTree(const vector<Info> &init_) {
        init(init_);
    }
    void init(int n_, Info v_ = Info()) {
        init(vector(n_, v_)); // seg(n + 1)
    }
    void init(const vector<Info> init_) {
        n = init_.size() - 1; // init_ 大小为 n + 1
        info.assign(4 << __lg(n), Info());
        tag.assign(4 << __lg(n), Tag());
        function<void(int, int, int)> build = [&](int p, int l, int r) {
            if (l == r) {
                info[p] = init_[l];
                info[p].l = l;
                info[p].r = r;

                return;
            }
            int m = l + r >> 1;
            build(p << 1, l, m);
            build(p << 1 | 1, m + 1, r);
            pull(p);
        };
        build(1, 1, n);
    }
    void pull(int p) {
        info[p] = info[p << 1] + info[p << 1 | 1];
    }
    void apply(int p, int l, int r, const Tag &v) {
        info[p].apply(v);
        tag[p].apply(v);
    }
    void push(int p, int l, int r) {
        int m = l + r >> 1;
        apply(p << 1, l, m, tag[p]);
        apply(p << 1 | 1, m + 1, r, tag[p]);
        tag[p] = Tag();
    }
    Info rangeQuery(int p, int l, int r, int x, int y) {
        if (x <= l && r <= y) {
            return info[p];
        }
        push(p, l, r);
        int m = l + r >> 1;
        if (y <= m) {
            return rangeQuery(p << 1, l, m, x, y);
        } else if (x >= m + 1) {
            return rangeQuery(p << 1 | 1, m + 1, r, x, y);
        } else {
            return rangeQuery(p << 1, l, m, x, y) + rangeQuery(p << 1 | 1, m + 1, r, x, y);
        }
    }
    Info rangeQuery(int l, int r) {
        return rangeQuery(1, 1, n, l, r);
    }
    void rangeModify(int p, int l, int r, int x, int y, const Tag &v) {
        if (x <= l && r <= y) {
            apply(p, l, r, v);
            return;
        }
        push(p, l, r);
        int m = l + r >> 1;
        if (y <= m) {
            rangeModify(p << 1, l, m, x, y, v);
        } else if (x >= m + 1) {
            rangeModify(p << 1 | 1, m + 1, r, x, y, v);
        } else {
            rangeModify(p << 1, l, m, x, y, v);
            rangeModify(p << 1 | 1, m + 1, r, x, y, v);
        }
        pull(p);
    }
    void rangeModify(int l, int r, const Tag &v) {
        rangeModify(1, 1, n, l, r, v);
    }
    void modify(int p, int l, int r, int x, const Info &v) {
        if (l == r) {
            info[p] = v;
            return;
        }
        push(p, l, r);
        int m = l + r >> 1;
        if (x <= m) {
            modify(p << 1, l, m, x, v);
        } else {
            modify(p << 1 | 1, m + 1, r, x, v);
        }
        pull(p);
    }
    void modify(int x, const Info &v) {
        modify(1, 1, n, x, v);
    }    
    template<class F>
    int findFirst(int p, int l, int r, int x, int y, F &&pred) {
        if (y < l || r < x) {
            return -1;
        }
        if (x <= l && r <= y && !pred(info[p])) {
            return -1;
        }
        if (l == r) {
            return l;
        }
        push(p, l, r);
        int m = l + r >> 1;
        int res = findFirst(p << 1, l, m, x, y, pred);
        if (res == -1) {
            res = findFirst(p << 1 | 1, m + 1, r, x, y, pred);
        }
        return res;
    }
    template<class F>
    int findFirst(int l, int r, F &&pred) {
        return findFirst(1, 1, n, l, r, pred);
    }
    template<class F>
    int findLast(int p, int l, int r, int x, int y, F &&pred) {
        if (y < l || r < x) {
            return -1;
        }
        if (x <= l && r <= y && !pred(info[p])) {
            return -1;
        }
        if (l == r) {
            return l;
        }
        push(p, l, r);
        int m = l + r >> 1;
        int res = findLast(p << 1 | 1, m + 1, r, x, y, pred);
        if (res == -1) {
            res = findLast(p << 1, l, m, x, y, pred);
        }
        return res;
    }
    template<class F>
    int findLast(int l, int r, F &&pred) {
        return findLast(1, 1, n, l, r, pred);
    }
}
// by lyh
\end{lstlisting}

\subsection{可持久化线段树}

维护历史区间最大值

\begin{lstlisting}[language=C++]
struct node{
    int lc, rc, dat;
}tr[M];

int build(int l, int r){
    int p = ++tot;
    if(l == r){
        tr[p].dat = a[l];
        return p;
    }
    int mid = (l+r)>>1;
    tr[p].lc = build(l, mid);
    tr[p].rc = build(mid+1, r);
    tr[p].dat = max(tr[tr[p].lc].dat, tr[tr[p].rc].dat);
    return p;
}

int insert(int k, int l, int r, int pos, int val){
    int p = ++tot;
    tr[p] = tr[k];
    if(l == r){
        tr[p].dat = val;
        return p;
    }
    int mid = (l+r)>>1;
    if(mid >= pos) tr[p].lc = insert(tr[k].lc, l, mid, pos, val);
    else tr[p].rc = insert(tr[k].rc, mid+1, r, pos, val);
    tr[p].dat = max(tr[tr[p].lc].dat, tr[tr[p].rc].dat);
    return p;
}

int query(int k, int l, int r, int L, int R){
    int res = -2e9;
    if(l >= L && r <= R) return tr[k].dat;
    int mid = (l+r)>>1;
    if(mid >= L) res = max(res, query(tr[k].lc, l, mid, L, R));
    if(mid < R) res = max(res, query(tr[k].rc, mid+1, r, L, R));
    return res; 
}
\end{lstlisting}

查询区间第 $k$ 小值（静态）

\begin{lstlisting}[language=C++]
#include<bits/stdc++.h>
using namespace std;

const int N = 2e5+10;
int n, m, tot, num;
int a[N], b[N], c[N], root[N];

struct node{
    int lc, rc, dat;
}tr[N<<5];

int build(int l, int r){
    int p = ++tot;
    if(l == r){
        tr[p].dat = 0;
        return p;
    }
    int mid = (l+r)>>1;
    tr[p].lc = build(l, mid);
    tr[p].rc = build(mid+1, r);
    tr[p].dat = tr[tr[p].lc].dat+tr[tr[p].rc].dat;
    return p;
}

int change(int k, int l, int r, int pos, int val){
    int p = ++tot;
    tr[p] = tr[k];
    if(l == r){
        tr[p].dat += val;
        return p;
    }
    int mid = (l+r)>>1;
    if(mid >= pos) tr[p].lc = change(tr[k].lc, l, mid, pos, val);
    else tr[p].rc = change(tr[k].rc, mid+1, r, pos, val);
    tr[p].dat = tr[tr[p].lc].dat+tr[tr[p].rc].dat;
    return p;
}

int query(int i, int j, int l, int r, int k){
    if(l == r) return l;
    int mid = (l+r)>>1;
    int cntl = tr[tr[j].lc].dat-tr[tr[i].lc].dat;
    if(cntl >= k) return query(tr[i].lc, tr[j].lc, l, mid, k);
    else return query(tr[i].rc, tr[j].rc, mid+1, r, k-cntl);
}

int Q(int x) {return lower_bound(c+1, c+num+1, x)-c;}

void discrete(){
    for(int i = 1; i <= n; i++) b[i] = a[i];
    sort(b+1, b+n+1);
    b[0] = -1;
    num = 0;
    for(int i = 1; i <= n; i++)
        if(b[i] != b[i-1]) c[++num] = b[i];
    for(int i = 1; i <= n; i++)
        a[i] = Q(a[i]);
}

int main(){
    ios::sync_with_stdio(0);
    cin.tie(0);
    cin >> n >> m;
    for(int i = 1; i <= n; i++)
        cin >> a[i];
    discrete();
    root[0] = build(1, num);
    for(int i = 1; i <= n; i++)
        root[i] = change(root[i-1], 1, num, a[i], 1);
    for(int i = 1; i <= m; i++){
        int l, r, k;
        cin >> l >> r >> k;
        int id = query(root[l-1], root[r], 1, num, k);
        cout << c[id] << endl;
    }
    return 0;
}
\end{lstlisting}

\subsection{可持久化并查集}

\subsection{李超线段树}

\begin{lstlisting}[language=C++]
using ll = long long;
using pii = pair<int, int>;
using pdi = pair<double, int>;
const int N = 1e5+10, modx = 39989, mody = 1e9+7;
const double eps = 1e-9;
struct line{
    double k, b;
    line(double _k = 0, double _b = 0): k(_k), b(_b) {}
}l[N];
int n, tot, tr[N<<2];

line get_line(double x1, double y1, double x2, double y2){
    if(x1 == x2) return line(0, max(y1, y2));
    return line((y1-y2)/(x1-x2), (x1*y2-x2*y1)/(x1-x2));
}

double calc(int i, ll x){
    return l[i].k*x+l[i].b;
}

int cmp(double x, double y){
    if(x-y > eps) return 1;
    if(y-x > eps) return -1;
    return 0;
}

pdi pmax(pdi a, pdi b){
    if(cmp(a.first, b.first) == 1) return a;
    if(!cmp(a.first, b.first))
        return a.second < b.second ? a : b;
    return b;
}

void update(int k, int l, int r, int u){
    int &v = tr[k], mid = (l+r)>>1;
    int cmid = cmp(calc(u, mid), calc(v, mid));
    if(cmid == 1 || (!cmid && u < v)) swap(u, v);
    int cl = cmp(calc(u, l), calc(v, l)), rl = cmp(calc(u, r), calc(v, r));
    if(cl == 1 || (!cl && u < v)) update(k<<1, l, mid, u);
    if(rl == 1 || (!rl && u < v)) update((k<<1)|1, mid+1, r, u);
}

void insert(int k, int l, int r, int L, int R, int u){
    if(l >= L && r <= R){
        update(k, l, r, u);
        return ;
    }
    int mid = (l+r)>>1;
    if(mid >= L) insert(k<<1, l, mid, L, R, u);
    if(mid < R) insert((k<<1)|1, mid+1, r, L, R, u);
}

pdi query(int k, int l, int r, ll x){
    int mid = (l+r)>>1;
    pdi res = {calc(tr[k], x), tr[k]};
    if(l == r) return res;
    if(x <= mid) res = pmax(res, query(k<<1, l, mid, x));
    if(x > mid) res = pmax(res, query((k<<1)|1, mid+1, r, x)); 
    return res;
}
\end{lstlisting}


\subsection{树状数组求第 $k$ 大}

\begin{lstlisting}[language=C++]

template<class T>
struct Fenwick {
    int n;
    vector<T> f;
    Fenwick(int n) {
        init(n);
    }
    void init(int n) {
        this->n = n;
        f.assign(n, 0);
    }
    void add(int x, T k) {
        assert(x > 0);
        for (; x < n; x += x & -x) {
            f[x] += k;
        }
    }
    T pre(int x) {
        assert(x < n);
        T res {};
        for (; x; x -= x & -x) {
            res += f[x];
        }
        return res;
    }
    T rangeSum(int l, int r) {
        return pre(r) - pre(l - 1);
    }
    int kth(int k) {
        int x = 0;
        T cur {};
        for (int i = 1 << __lg(n); i; i >>= 1) {
            if (x + i < n && cur + f[x + i] < k) {
                cur += f[x + i];
                x += i;
            }
        }
        return x + 1;
    }
};
\end{lstlisting}

\subsection{笛卡尔树}

每个节点为一对键值 $(k, w)$，其中 $k$ 满足二叉搜索树的性质，$w$ 满足堆的性质，通过单调栈维护右链建树，新的节点插入到右链，将 $w$ 大于该节点的子树作为左子树。

\begin{lstlisting}[language = C++]
// stk 维护笛卡尔树中节点对应到序列中的下标
for (int i = 1; i <= n; i++) {
    int k = top;  // top 表示操作前的栈顶，k 表示当前栈顶
    while(k > 0 && w[stk[k]] > w[i]) k--;  // 维护右链上的节点
    if(k) rs[stk[k]] = i;  // 栈顶元素.右儿子 := 当前元素
    if(k < top) ls[i] = stk[k + 1]; 
    // 当前元素.左儿子 := 上一个被弹出的元素  
    stk[++k] = i; // 当前元素入栈
    top = k;
}
\end{lstlisting}

\subsection{pbds tree}

\begin{lstlisting}[language=C++]
#include <bits/extc++.h>
using namespace __gnu_pbds;
using pbds_tree = tree<array<int, 2>, null_type, less<>, rb_tree_tag, tree_order_statistics_node_update>;

void solve() {
    int idx = 0; // 有重复元素，要设置唯一索引
    pbds_tree t;
    while (q--) {
        if (o == 1) {
            // 插入
            t.insert({x, idx++});
        } else if (o == 2) {
            // 删除
            auto it = t.lower_bound({x, 0});
            if (it != t.end() && (*it)[0] == x) {
                t.erase(it);
            }
        } else if (o == 3) {
            // 排名
            cout << t.order_of_key({x, 0}) + 1 << "\n"; 
        } else if (o == 4) {
            // 第 x 大，find_by_order 是 0-based，要查 x - 1
            cout << (*t.find_by_order(x - 1))[0] << "\n";
        } else if (o == 5) {
            // 前驱
            cout << (*--t.lower_bound({x, 0}))[0] << "\n";
        } else {
            // 后继
            cout << (*t.upper_bound({x, idx}))[0] << "\n";
        }
    }
}
\end{lstlisting}

\subsection{pbds pq}

\begin{lstlisting}[language=C++]
using pbds_pq = __gnu_pbds::priority_queue<array<int, 2>, greater<>>;

void solve() {
    int n, m;
    cin >> n >> m;
    
    vector<pbds_pq> pq(n + 1);
    for (int i = 1; i <= n; i++) {
        int x;
        cin >> x;
        pq[i].push({x, i});
    }
    
    vector<bool> del(n + 1);
    DSU dsu(n + 1);

    while (m--) {
        int o;
        cin >> o;
        
        if (o == 1) {
            // 合并第 x 个数和第 y 个数所在的堆
            int x, y;
            cin >> x >> y;
            
            int fx = dsu.find(x);
            int fy = dsu.find(y);
            if (del[x] || del[y] || fx == fy) {
                continue;
            }
            if (pq[fx].size() < pq[fy].size()) {
                swap(fx, fy);
            }
            pq[fx].join(pq[fy]);
            dsu.merge(fx, fy);

        } else {
            // 删除第 x 个数所在堆的堆顶
            int x;
            cin >> x;

            if (del[x]) {
                cout << "-1\n";
            } else {
                x = dsu.find(x);

                auto [val, id] = pq[x].top();
                pq[x].pop();

                del[id] = true;
                cout << val << "\n";
            }
        }
    }
}
\end{lstlisting}

\subsection{树状数组套线段树}

\begin{lstlisting}[language=C++]
int Y = 2e5;
struct Node {
    int l = 0, r = 0;
    int ls = 0, rs = 0;
    int sum = 0, add = 0;
};
vector<Node> in;
int newNode(int l, int r) {
    Node node;
    node.l = l;
    node.r = r;
    in.push_back(node);
    return in.size() - 1;
}
void push(int p) {
    if (in[p].add == 0 || in[p].l == in[p].r) {
        return;
    }
    int m = (in[p].l + in[p].r) >> 1;
    if (!in[p].ls) {
        in[p].ls = newNode(in[p].l, m);
    }
    if (!in[p].rs) {
        in[p].rs = newNode(m + 1, in[p].r);
    }
    in[in[p].ls].sum += (m - in[p].l + 1) * in[p].add;
    in[in[p].ls].add += in[p].add;
    in[in[p].rs].sum += (in[p].r - m) * in[p].add;
    in[in[p].rs].add += in[p].add;
    in[p].add = 0;
}
void pull(int p) {
    in[p].sum = 0;
    if (in[p].ls) {
        in[p].sum += in[in[p].ls].sum;
    }
    if (in[p].rs) {
        in[p].sum += in[in[p].rs].sum;
    }
}
int rangeModify(int p, int l, int r, int yl, int yr, int k) {
    if (yr < l || r < yl) {
        return p;
    }
    if (!p) {
        p = newNode(l, r);
    }
    if (yl <= l && r <= yr) {
        in[p].sum += (r - l + 1) * k;
        in[p].add += k;
        return p;
    }
    push(p);
    int m = (l + r) >> 1;
    in[p].ls = rangeModify(in[p].ls, l, m, yl, yr, k);
    in[p].rs = rangeModify(in[p].rs, m + 1, r, yl, yr, k);
    pull(p);
    return p;
}
int rangeQuery(int p, int yl, int yr) {
    if (!p || yr < in[p].l || in[p].r < yl) {
        return 0;
    }
    if (yl <= in[p].l && in[p].r <= yr) {
        return in[p].sum;
    }
    push(p);
    return rangeQuery(in[p].ls, yl, yr) + rangeQuery(in[p].rs, yl, yr);
}

int X = 2e5;
struct OutNode {
    int inRoot = 0;
};
vector<OutNode> out;
void add(int x, int yl, int yr, int k) {
    assert(x > 0);
    for (; x <= X; x += x & -x) {
        out[x].inRoot = rangeModify(out[x].inRoot, 1, Y, yl, yr, k);
    }
}
int pre(int x, int yl, int yr) {
    assert(x <= X);
    int res = 0;
    for (; x; x -= x & -x) {
        res += rangeQuery(out[x].inRoot, yl, yr);
    }
    return res;
}
\end{lstlisting}

\subsection{ST 表}

\begin{lstlisting}[language=C++]
template<class T>
struct SparseTable {
    using Func = function<T(T, T)>; // lambda
    int n;
    vector<int> log2;
    vector<vector<T>> st;
    Func func;

    SparseTable(const vector<T> &a, Func f) {
        init(a, f);
    }
    void init(const vector<T> &a, Func f) {
        n = a.size();
        func = f;
        int k = __lg(n) + 1;
        log2.assign(n, 0);
        st.assign(n, vector<T>(k));
        for (int i = 2; i < n; i++) {
            log2[i] = log2[i >> 1] + 1;
        }
        for (int i = 1; i < n; i++) {
            st[i][0] = a[i];
        }
        for (int j = 1; j < k; j++) {
            for (int i = 1; i + (1 << j) - 1 < n; i++) {
                st[i][j] = func(st[i][j - 1], st[i + (1 << j - 1)][j - 1]);
            }
        }
    }
    T query(int l, int r) {
        int k = log2[r - l + 1];
        return func(st[l][k], st[r - (1 << k) + 1][k]);
    }
};
// by lyh
\end{lstlisting}

\subsection{XOR Hash}

\begin{lstlisting}[language=C++]
mt19937_64 rng(chrono::steady_clock::now().time_since_epoch().count());
struct XorHash {
    int n;
    vector<int> Base, Xor, Hash;
    XorHash(const vector<int> &a) {
        init(a);
    }
    void init(const vector<int> &a) {
        n = a.size();
        Base.assign(n, 0);
        Xor.assign(n, 0);
        Hash.assign(n, 0);
        for (int i = 1; i < n; i++) {
            Base[i] = rng();
            Xor[i] = Xor[i - 1] ^ Base[i];
        }
        for (int i = 1; i < n; i++) {
            Hash[i] = Hash[i - 1] ^ Base[a[i]];
        }
    }
    bool isPemutation(int l, int r, int s, int t) { // a[l, r]是否是整数[s, t]的排列
        return (Hash[r] ^ Hash[l - 1]) == (Xor[t] ^ Xor[s - 1]);
    }
};
// by lyh
\end{lstlisting}

\subsection{动态图连通性 (线段树分治+可撤销并查集)}

利用线段树对询问进行划分, 每个叶子节点对应一个时间点. 记每条边的存在时间为 $[l, r]$, 将该区间插入线段树中, 线段树的每个节点记录对应时间段内存在的边. 建完树后 dfs 遍历线段树, 访问至叶节点回答询问. 回溯时撤销当前节点对并查集的操作.
\begin{lstlisting}[language=C++]
// LOJ 121
#include<bits/stdc++.h>
using namespace std;

using pii = pair<int, int>;
const int N = 5e3+10, M = 5e5+10;
map<pii, int> mp;
int n, m, ans[M];
pii q[M];
vector<pii> tr[M*20];

struct operation{
    int x, y, s_x;
};
struct DSU{
    int num;
    vector<int> f, s;
    stack<operation> op;
    DSU() {}
    DSU(int sz){
        num = sz;
        f = vector<int> (num+1, 0);
        s = vector<int> (num+1, 1);
        iota(f.begin(), f.end(), 0);
    }
    int find(int x){
        if(f[x] == x) return x;
        return find(f[x]);
    }
    void merge(int x, int y){
        x = find(x); y = find(y);
        if(s[x] > s[y]) swap(x, y);
        op.push({x, y, s[x]});
        if(x != y){
            f[x] = y;
            s[y] += x; s[x] = 0;
        }
    }
    void split(){
        auto [x, y, s_x] = op.top();  
        op.pop();
        if(x == y) return ;
        f[x] = x;
        s[x] = s_x; s[y] -= s_x;
    }
    void back(int x){
        for(int i = 1; i <= x; i++) split();
    }
}g;
void insert(int k, int l, int r, int L, int R, pii e){
    if(l >= L && r <= R){
        tr[k].push_back(e);
        return ;
    }
    int mid = (l+r)>>1;
    if(mid >= L) insert(k<<1, l, mid, L, R, e);
    if(mid < R) insert((k<<1)|1, mid+1, r, L, R, e);
}
void dfs(int k, int l, int r){
    for(auto [x, y]: tr[k])
        g.merge(x, y);
    if(l == r){
        auto [x, y] = q[l];
        if(x == 0 && y == 0) {g.back(tr[k].size()); ans[l] = -1; return ;}
        if(g.find(x) == g.find(y)) ans[l] = 1;
        else ans[l] = 0;
        g.back(tr[k].size());
        return ;
    }
    int mid = (l+r)>>1;
    dfs(k<<1, l, mid);
    dfs((k<<1)|1, mid+1, r);
    g.back(tr[k].size());
}

int main(){
    ios::sync_with_stdio(0);
    cin >> n >> m;
    for(int i = 1; i <= m; i++){
        int op, x, y;
        cin >> op >> x >> y;
        if(x > y) swap(x, y);
        q[i] = {0, 0};
        if(op == 0) mp[{x, y}] = i;
        else if(op == 1){
            insert(1, 1, m, mp[{x, y}], i-1, {x, y});
            mp[{x, y}] = 0;
        }
        else q[i] = {x, y};
    }
    for(auto [e, i]: mp)
        if(i) insert(1, 1, m, i, m, e);
    g = DSU(n);
    dfs(1, 1, m);
    for(int i = 1; i <= m; i++) if(ans[i] != -1){
        if(ans[i]) cout << "Y\n";
        else cout << "N\n";
    }
    return 0;
}
\end{lstlisting}